\section*{Bab \ref{ch:g12:subduction-magmatism} --- Magmatisme di Zona Subduksi}

\begin{solution}{exo:g12:subduction-magmatism:1}
(a) \omterm{def:g12:subduction-magmatism:arc}{busur gunung api}; (b) \omterm{def:g12:subduction-magmatism:hydrous-melting}{pelelehan hidrus} (pelelehan fluks); (c) \omterm{def:g12:subduction-magmatism:granodiorite}{granodiorit};
(d) \omterm{def:g12:subduction-magmatism:batholith}{batolit}; (e) \omterm{def:g12:subduction-magmatism:continental-growth}{pertumbuhan benua}.
\end{solution}

\begin{solution}{exo:g12:subduction-magmatism:2}
(1) Setidaknya $24/0.25 = \qty{96}{km/h}$, yaitu $\num{96000}/3600 \approx
\qty{27}{m/s}$. (2) Semburan itu setidaknya $480/96 = 5$ kali lebih cepat.
\end{solution}

\begin{solution}{exo:g12:subduction-magmatism:3}
(1) $838/1214 \approx 0.690$, sekitar \qty{69}{\percent}. (2) Di sisi lempeng
yang berada di atas, yang tidak tenggelam. (3) Di antara palung dan front,
\omterm{def:g11:subduction-zones:slab}{lempengan} belum cukup dalam: \omterm{def:g11:subduction-zones:slab}{lempengan} itu belum mencapai kedalaman sekitar
seratus kilometer, tempat airnya dilepaskan ke dalam mantel yang cukup panas
untuk meleleh setelah basah.
\end{solution}

\begin{solution}{exo:g12:subduction-magmatism:4}
(1) $45/(45 + 15) = 0.75$: \qty{75}{\percent} feldsparnya berupa plagioklas.
(2) Ya: batuan itu berbutir kasar, mengandung lebih dari \qty{20}{\percent}
kuarsa, dan bagian plagioklasnya terletak antara 65 dan
\qty{90}{\percent}.
\end{solution}

\begin{solution}{exo:g12:subduction-magmatism:5}
Magma busur lebih kaya silika, sehingga lebih kental; dan magma itu membawa air
yang membentuknya, sehingga kaya gas. Kekentalan yang tinggi dan kandungan gas
yang tinggi sama-sama mempermudah letusan eksplosif.
\end{solution}

\begin{solution}{exo:g12:subduction-magmatism:6}
(1) $H = 150 \times \tan\ang{35} \approx 150 \times 0.700 \approx
\qty{105}{km}$. (2) Itulah rata-rata dunia.
\end{solution}

\begin{solution}{exo:g12:subduction-magmatism:7}
$d = H/\tan\theta$: $105/\tan\ang{20} \approx 105/0.364 \approx
\qty{288}{km}$ pada \ang{20}, dan $105/\tan\ang{45} = \qty{105}{km}$ pada
\ang{45}. Makin landai \omterm{def:g11:subduction-zones:slab}{lempengan}, makin jauh busur dari palung.
\end{solution}

\begin{solution}{exo:g12:subduction-magmatism:8}
(1) $T = -5.104 \times 3.16^2 + 132.9 \times 3.16 + 1120.7 \approx -51.0 +
420.0 + 1120.7 \approx \qty{1490}{\celsius}$. (2) Sekitar $1490 - 1000 =
\qty{490}{\celsius}$. (3) Kering, tidak: \qty{1300}{\celsius} berada di bawah
\qty{1490}{\celsius}. Basah, ya: suhu itu di atas \qty{1000}{\celsius}.
\end{solution}

\begin{solution}{exo:g12:subduction-magmatism:9}
(1) $2.5/2.7 \approx 0.93$ tahun, sekitar 340 hari. (2) Longsoran itu
memindahkan batuan tua yang sudah ada; busur membuat batuan baru dari mantel.
Hanya volumenya yang dibandingkan.
\end{solution}

\begin{solution}{exo:g12:subduction-magmatism:10}
(1) $2.7 \times \num{1000000} = \qty{2.7e6}{km^3}$. (2) $\num{9e9}/2.7
\approx \num{3.3e9}$ tahun, 3.3 miliar tahun. (3) Waktu itu kurang dari umur
Bumi, sehingga busur dapat saja membuat seluruh \omterm{def:g12:continental-crust:continental-crust}{kerak benua}; tetapi \omterm{def:g11:subduction-zones:subduction}{subduksi}
juga membawa kembali kerak, sehingga sebenarnya sebagian besar kerak telah
dibuat dan dihancurkan lebih dari sekali.
\end{solution}

\begin{solution}{exo:g12:subduction-magmatism:11}
(1) Gempa bumi dangkal juga terjadi di lempeng yang berada di atas dan pada
sesar di antara kedua lempeng; hanya gempa-gempa dalam yang terletak di dalam
\omterm{def:g11:subduction-zones:slab}{lempengan} saja.
(2) $116 - 80 = \qty{36}{km}$ pada jarak \qty{86.4}{km}: $\tan\theta = 36/86.4
\approx 0.417$, $\theta \approx \ang{22.6}$.
\end{solution}

\begin{solution}{exo:g12:subduction-magmatism:12}
Dalam segitiga siku-siku proposisi itu, \omterm{def:g11:subduction-zones:slab}{lempengan} adalah sisi miring $L$ dan
$d$ sisi yang berdampingan dengan sudut $\theta$, sehingga $\cos\theta = d/L$ dan $L =
d/\cos\theta$. Untuk Jepang, $L = 270/\cos\ang{20} \approx 270/0.940 \approx
\qty{287}{km}$. Dengan laju \qty{83}{mm} setahun: $\num{2.87e8}/83 \approx \num{3.5e6}$
tahun, sekitar 3.5 juta tahun.
\end{solution}

\begin{solution}{exo:g12:subduction-magmatism:13}
(1) $V = \pi r^2 h = \pi \times 0.5^2 \times 8 \approx \qty{6.3}{km^3}$.
(2) Sekitar $6.3/2.5 \approx 2.5$ kali volume longsoran itu.
(3) $P = \rho g z$: pada \qty{6}{km}, $2830 \times 9.81 \times 6000 \approx
\qty{1.67e8}{Pa} = \qty{0.17}{GPa}$; pada \qty{14}{km}, sekitar
\qty{3.89e8}{Pa} = \qty{0.39}{GPa}.
\end{solution}

\begin{solution}{exo:g12:subduction-magmatism:14}
(1) $2.7 - 3.2 = \qty{-0.5}{km^3}$ setahun, kehilangan \qty{5e5}{km^3} per
sejuta tahun. (2) $500 \times \num{5e5} = \qty{2.5e8}{km^3}$, yaitu
$\num{2.5e8}/\num{9e9} \approx \qty{2.8}{\percent}$ dari kerak. (3) Setiap
laju adalah perkiraan, dengan ketidakpastian yang jauh lebih besar daripada
selisih \qty{0.5}{km^3} setahun di antara keduanya: dalam batas galatnya,
keduanya seimbang.
\end{solution}

\begin{solution}{exo:g12:subduction-magmatism:15}
(1) $\num{20000} \times 10 = \qty{2e5}{km^3}$. (2) $\num{2e5}/1 =
\num{200000}$ tahun untuk seluruh busur benua di dunia bersama-sama. Satu busur
hanya menerima sebagian kecil magma itu, sehingga satu \omterm{def:g12:subduction-magmatism:batholith}{batolit} memerlukan waktu
jauh lebih lama: \omterm{def:g12:subduction-magmatism:batholith}{batolit} dibangun \omterm{def:g11:igneous-rocks:intrusion}{pluton} demi \omterm{def:g11:igneous-rocks:intrusion}{pluton}, selama jutaan tahun.
(3) Setidaknya batuan setebal \qty{5}{km} di atas bagian puncaknya.
\end{solution}

\begin{solution}{pb:g12:subduction-magmatism:1}
\textbf{1.} $111.2 \times \cos\ang{39} \approx 111.2 \times 0.777 \approx
\qty{86.4}{km}$.

\textbf{2.} $(144.1 - 141.0) \times 86.4 \approx 3.1 \times 86.4 \approx
\qty{268}{km}$, sekitar \qty{270}{km}.

\textbf{3.} $(141.0 - 139.7) \times 86.4 \approx \qty{112}{km}$.

\textbf{4.} Di atas lempeng yang berada di atas, yang tidak tenggelam, di barat
palung: Lempeng Pasifik tenggelam di bawahnya.

\textbf{5.} Di sana \omterm{def:g11:subduction-zones:slab}{lempengan} belum cukup dalam untuk melepaskan airnya ke
dalam mantel yang cukup panas untuk meleleh.

\textbf{6.} $116 - 80 = \qty{36}{km}$.

\textbf{7.} Satu derajat bujur, \qty{86.4}{km}.

\textbf{8.} $\tan\theta = 36/86.4 \approx 0.417$, sehingga $\theta \approx
\ang{22.6}$.

\textbf{9.} Garis itu naik \qty{36}{km} setiap \qty{86.4}{km} ke arah timur,
sehingga mencapai permukaan $98/0.417 \approx \qty{235}{km}$ di timur Iwate, pada
$141.0 + 235/86.4 \approx \ang{143.7}$~BT: sekitar $(144.1 - 143.7) \times 86.4
\approx \qty{35}{km}$ di barat palung.

\textbf{10.} \omterm{def:g11:subduction-zones:slab}{Lempengan} itu bukan bidang datar: \omterm{def:g11:subduction-zones:slab}{lempengan} memasuki palung dalam
keadaan hampir mendatar dan makin melengkung ke bawah selagi tenggelam.

\textbf{11.} $268 \times \tan\ang{20} \approx 268 \times 0.364 \approx
\qty{98}{km}$.

\textbf{12.} Keduanya sesuai: \qty{98}{km} dengan kedua cara.

\textbf{13.} Keduanya terletak di dalam rentang dunia, sedikit lebih dangkal
daripada rata-rata \qty{105}{km}.

\textbf{14.} Megata terletak \qty{112}{km} di barat Iwate, sehingga \omterm{def:g11:subduction-zones:slab}{lempengan}
di bawahnya sedalam sekitar $98 + 0.417 \times 112 \approx \qty{145}{km}$.

\textbf{15.} Di dalam mantel, sekitar $98 - 35 = \qty{63}{km}$ di bawah dasar
kerak.

\textbf{16.} Air itu masuk di dekat punggungan tempat kerak terbentuk, dari air
laut; sebagian besar air itu tersimpan di dalam serpentin.

\textbf{17.} $T = -5.104 \times 3.16^2 + 132.9 \times 3.16 + 1120.7 \approx
\qty{1490}{\celsius}$.

\textbf{18.} Sebanyak $1490 - 1000 = 490$ sampai $1490 - 950 = \qty{540}{\celsius}$:
sekitar \qty{500}{\celsius}.

\textbf{19.} \omterm{def:g11:subduction-zones:slab}{Lempengan} itu beberapa ratus derajat lebih dingin daripada mantel
di sekitarnya, terlalu dingin untuk meleleh bahkan dengan air: baji mantel
panas di atasnya, setelah dibasahi, itulah yang meleleh.

\textbf{20.} \omterm{def:g11:subduction-zones:slab}{Lempengan} itu terletak sekitar \textbf{\qty{100}{km} di bawah front
vulkanik} Jepang timur laut (\qty{98}{km} dengan kedua cara).
\end{solution}
